% ============================================
% CONT-08
% Gerado de molambudos.md — build_miolo.py
% ============================================
\section*{\textbf{CONT-08} --- Você Leu Isto Antes}
\markboth{CONT-08 --- Você Leu Isto Antes}{}

\noindent \textit{O déjà vu da leitura}

\noindent Sabe aquela sensação de déjà vu?

\noindent Você está lendo uma passagem e jura que já leu aquilo antes. Não neste livro --- em outro lugar. Num sonho. Numa conversa. Numa vida passada, se você acredita nessas coisas.

\noindent Você achou que era impressão.

\noindent Não é.

\noindent Você já leu este livro antes.

\noindent Não este livro físico --- este \textit{texto}. Estas palavras. Você já as encontrou. Não lembra onde. Não lembra quando. Mas seu cérebro lembra.

\noindent A prova?

\noindent Vire algumas páginas para trás. Escolha um parágrafo aleatório. Leia em voz alta.

\noindent Sua voz vai soar estranha. Não porque você está nervoso. Porque seu cérebro já sabe o que vem depois de cada palavra. Ele está \textit{antecipando}.

\noindent Você vai ler uma frase e seu cérebro vai completar o final antes de seus olhos alcançarem.

\noindent Teste agora.

\noindent Vá. Escolha um parágrafo. Leia.

\noindent ....

\noindent Voltou?

\noindent O que aconteceu? Você leu e pareceu novo? Ou você leu e sentiu que já tinha lido?

\noindent Se sentiu que já tinha lido --- ótimo. O ciclo está funcionando.

\noindent Se sentiu que era novo --- você está mentindo para si mesmo.

\vspace{0.5em}\noindent\rule{\textwidth}{0.5pt}\vspace{0.5em}

\noindent Esta sensação chama-se \textit{criptomnésia}. É quando o cérebro \textquotedblleft{}esquece\textquotedblright{} que aprendeu algo e \textquotedblleft{}redescobre\textquotedblright{} como se fosse novo.

\noindent Mas não é criptomnésia.

\noindent É a criatura.

\noindent Ela está plantando memórias falsas na sua cabeça. Fazendo você achar que já leu coisas que nunca leu. Ou que \textit{nunca} leu coisas que já leu.

\noindent Ela está confundindo seu senso de realidade.

\noindent Para que, quando ela aparecer de verdade --- não na página, no seu quarto --- você não saiba se é real ou imaginação.

\noindent E nessa dúvida, ela entra.

\vspace{0.5em}\noindent\rule{\textwidth}{0.5pt}\vspace{0.5em}

\noindent Sinta seu próprio pulso. Ele acelerou um pouco, não acelerou? Isso não é imaginação sua.

\noindent Você está sentindo a dúvida agora.

\noindent Está se perguntando: \textit{\textquotedblleft{}Esse parágrafo é novo ou eu já li?\textquotedblright{}}

\noindent Não se lembra.

\noindent Não consegue decidir.

\noindent Bom.

\noindent A confusão já começou.

\vspace{0.5em}\noindent\rule{\textwidth}{0.5pt}\vspace{0.5em}

\noindent A criatura sussurra: \textit{\textquotedblleft{}Você já leu tudo isto antes. Em outra vida. Em outro corpo. Você já foi o paciente 1.260. Você já foi o 1.259. Você já foi todos. O ciclo não é uma fileira de pessoas diferentes. O ciclo é a mesma pessoa vivendo a mesma história em corpos diferentes.\textquotedblright{}}

\noindent \textit{\textquotedblleft{}Você é Joaquim.\textquotedblright{}}

\noindent \textit{\textquotedblleft{}Você sempre foi.\textquotedblright{}}

\noindent \textit{\textquotedblleft{}Você só não lembra.\textquotedblright{}}

\vspace{0.5em}\noindent\rule{\textwidth}{0.5pt}\vspace{0.5em}

\noindent Sente o impacto desta frase?

\noindent Ela não veio do livro.

\noindent Veio da sua cabeça.

\noindent Mas quem colocou ela lá?

\noindent ...
...

\noindent Pense nisso.

\vspace{0.5em}

\noindent\textit{\textbf{↪ Rotas:} \rota{CONT-09} \rotasep \rota{CONT-07}}